% !TeX root = ../main.tex

\chapter{結論與未來工作}
\label{chap:conclusion}

\section{研究總結}
\label{sec:summary}

本論文針對自 PGP 9.5 染色之皮膚切片影像自動量測表皮內神經纖維密度 (IENFD) 之實務問題，指出既有語意分割模型之輸出常因染色不均而呈現\textbf{破碎化}，單純骨架化並無法得到符合 EFNS 規則之可計數拓樸。為解決此片段重組問題，本研究將其形式化為基於\emph{像素成本圖}上之多源最短路徑與圖論重建，並提出 \textbf{標註擴張連結演算法} (Annotation-Grow Linker, AGL)。AGL 以綠色通道經形態學去背景、CLAHE 與 Sato 血管性濾波而得之成本圖為基底，將每個標註連通元件視為獨立來源同步執行 Dijkstra 擴張；於擴張場中搜尋相鄰元件之最小成本相遇點以建立元件層圖，並以成本閾值裁剪後求最小生成森林，最後將回溯路徑與原標註合併、骨架化以產生可供分段偵測、區域標記與有效跨越計數之像素級拓樸圖。

相較於既有基準方法，AGL 具備以下三項主要優勢：

\begin{itemize}
  \item \textbf{影像感知之連結準則}：連結路徑自動遵循影像上實際存在之纖維線索，避免純距離連結於高密度區域產生橫向錯接。
  \item \textbf{兩階段拓樸過濾}：閾值裁剪抑制偽連線、MST 抑制冗餘連線，兩者皆於元件層上操作，計算成本低且可解釋。
  \item \textbf{端到端可稽核流程}：輸出為可視化之拓樸圖而非純量結果，每一次被計為跨越之片段皆對應影像上具體之像素路徑，便於臨床複核。
\end{itemize}

於本研究所使用之資料集上，AGL 在跨越計數之絕對誤差與 Pearson 相關性上均優於 Pure-MST、DBSCAN-Linker 與 Raw-Skeleton；消融實驗亦顯示指數成本函數、僅向下之 ROI 膨脹、以及中等範圍之裁剪閾值 $\tau$ 三者皆對最終結果有實質貢獻。

\section{研究限制}
\label{sec:limitations}

本研究之結果受限於以下若干因素：

\begin{enumerate}
  \item \textbf{資料集規模與多樣性}。本文所使用之資料集來源單一，切片厚度與染色協定一致，其結論向其他醫院或切片規格之遷移性仍待驗證。
  \item \textbf{對上游分割品質之依賴}。AGL 僅能以「既有標註」為擴張來源，若某纖維於整段長度上皆未被語意分割模型捕捉，AGL 將無法自無中生有地加入該纖維。
  \item \textbf{表皮遮罩之精確性}。跨越事件之判定直接依賴 $M$，在手工標註誤差或自動表皮分割失準時會造成計數偏移。
  \item \textbf{二維切片之限制}。本文所處理之影像為單一垂直切片之二維投影；真實纖維於三維空間中可能繞出又進入同一切片，於二維上被計為多次跨越。此問題需透過真正的三維影像處理方能徹底解決。
\end{enumerate}

\section{未來工作}
\label{sec:future}

基於上述結果與限制，本文提出以下幾項具體之延伸方向：

\subsection{與語意分割之聯合訓練}

現行 AGL 將語意分割之輸出視為固定輸入。若將 AGL 之重建結果反饋至分割模型之訓練目標，使分割模型偏好輸出「易於被 AGL 補足之片段」而非「形狀完整但本質錯誤之大塊」，有機會形成一個\textbf{分割—重建雙向優化}之迴圈。一種具體之實作是：將 AGL 輸出之重建骨架與專家跨越地面真值之差異回傳為分割訓練之 perception loss。

\subsection{三維切片之延伸}

隨著組織透明化 (tissue clearing) 與共聚焦／螢光三維掃描技術之成熟，研究級資料已可取得三維之神經纖維影像。AGL 之架構自然可延伸至三維：僅需將 8-連通鄰域改為 26-連通，並於 3D Hessian 矩陣上計算 Frangi vesselness 即可。在三維中，跨越 DEJ 之判定亦需改為與一個\emph{DEJ 曲面}之相交關係，而非與 2D 遮罩之相交；未來工作將探討此延伸之數值與效能問題。

\subsection{可學習之成本函數}

本文之成本函數為手動設計之固定映射。若以小型卷積網路將影像映射為成本圖、並以 AGL 之下游計數誤差為監督訊號進行反向傳播 (可藉由 soft Dijkstra 或 differentiable shortest path 實作)，則有機會進一步提升成本圖對 DEJ 附近染色中斷之適應性。

\subsection{每元件停止閾值之自適應}

目前每元件停止閾值 $\tau_i$ 以種子像素之成本統計量定義，但在極小元件上該統計量之變異較大。未來可改以\textbf{元件鄰域之空間先驗}（如同切片內之纖維走向先驗、鄰近元件之距離分布）作為動態停止條件，以降低在局部雜訊場景下之擴張範圍。

\subsection{不確定性量化}

對於臨床決策，除了 IENFD 之點估計外，其不確定性亦具重要意義。未來可考慮以 bootstrap 或以 Dijkstra 擴張之候選相遇點排名作為信心分數，對每一條被計為跨越之片段輸出對應之信心值，並進一步量化整張切片 IENFD 之信賴區間。

\section{結語}
\label{sec:closing}

自動化之 IENFD 量測是一項同時需要影像處理、圖論重建與臨床知識之整合性任務；單憑其中一環皆難以兼顧計數之準確性與輸出之可解釋性。本論文提出之 AGL 展示了\emph{將手工設計之影像先驗與圖論方法結合}之可行性，並於一個模組化之實作中維持每個階段之獨立可替換性。期望此研究能作為後續「分割—重建—計數」整合式研究之基礎，最終推動小纖維神經病變之客觀量測在臨床上之廣泛採用。
